% =====================================================================
%  IEEE 会议论文骨架（ICASSP / 通用双栏）
%  用法：把本文件当起点，替换 [占位内容]。需要官方 IEEEtran.cls。
%  编译：pdflatex -> bibtex -> pdflatex -> pdflatex
%  注意：投稿为双盲时，请勿在正文/元数据出现作者与机构信息。
% =====================================================================
\documentclass{article}
\usepackage{spconf,amsmath,graphicx}   % spconf 为 ICASSP 常用样式；无则改用 IEEEtran

% ---- 双盲：投稿版作者留空/匿名；Camera-ready 再补真实信息 ----
\title{[方法名]: [机制/亮点] for [任务]}
\name{Anonymous}                       % 双盲：Anonymous ; 终稿：Author A, Author B
\address{Anonymous Institution}        % 双盲：Anonymous ; 终稿：真实机构

\begin{document}
\maketitle

% =========================== Abstract ================================
\begin{abstract}
% 5 句骨架：背景/问题 -> 现有不足 -> 方法 -> 主要结果(具体数字) -> 意义
[一句话背景与痛点。] [现有方法的不足。] We propose [方法],
which [核心机制]. Experiments on [数据集] show [具体数字，如 X dB / Y\%]
improvement over the best baseline. [结论意义。]
\end{abstract}

\begin{keywords}
[关键词1], [关键词2], [关键词3], [关键词4]
\end{keywords}

% ========================= 1. Introduction ===========================
\section{Introduction}
\label{sec:intro}
% 背景与重要性
Recently, [任务] has attracted growing attention because [原因].
% 现有不足 / gap
However, existing approaches struggle to [问题], mainly because [原因].
% 本文做法 + 编号贡献
In this paper, we propose [方法] to address this issue.
The main contributions of this work are:
\begin{itemize}
  \item [贡献一：做了什么 + 带来什么]
  \item [贡献二]
  \item [贡献三]
\end{itemize}

% ========================= 2. Related Work ===========================
\section{Related Work}
\label{sec:related}
Prior studies can be broadly divided into three lines: [路线一], [路线二],
and [路线三]. Unlike [A] that focuses on [X], our method addresses [Y] by [Z].

% ========================= 3. Proposed Method ========================
\section{Proposed Method}
\label{sec:method}
% 3.1 问题形式化
\subsection{Problem Formulation}
Let $x \in \mathbb{R}^{T}$ denote the observed signal; our goal is to
estimate the target $s$ from $x$.
% 3.2 总览（建议放一张 block diagram）
\subsection{Overview}
\begin{figure}[t]
  \centering
  % \includegraphics[width=\linewidth]{fig/overview.pdf}
  \caption{Overview of the proposed [方法].}
  \label{fig:overview}
\end{figure}
% 3.3 逐模块 + 直觉解释
\subsection{[Key Module]}
Intuitively, this module lets the network focus on [作用]. Formally,
\begin{equation}
  y = f_{\theta}(x),
  \label{eq:main}
\end{equation}
where $\theta$ denotes the learnable parameters.

% ========================= 4. Experiments ===========================
\section{Experiments}
\label{sec:exp}
\subsection{Setup}
% 数据集 / 指标 / 基线 / 训练细节
We evaluate on [数据集] and report [指标]. We compare against [基线].

\subsection{Main Results}
\begin{table}[t]
  \centering
  \caption{Comparison with baselines. Best in \textbf{bold}.}
  \label{tab:main}
  \begin{tabular}{lccc}
    \hline
    Method & Metric1 & Metric2 & \#Params \\
    \hline
    Baseline A & 00.0 & 00.0 & 0.0M \\
    \textbf{Ours} & \textbf{00.0} & \textbf{00.0} & \textbf{0.0M} \\
    \hline
  \end{tabular}
\end{table}
As shown in Table~\ref{tab:main}, our method outperforms [基线] by
[X] in [指标] while using [Y] fewer parameters.

\subsection{Ablation Study}
% 逐个去掉模块，证明每部分有用
Table~\ref{tab:ablation} shows that removing [模块] degrades performance
by [数字], confirming its contribution.

% ========================= 5. Conclusion ============================
\section{Conclusion}
\label{sec:conclusion}
We presented [方法], demonstrating that [核心发现].
Future work includes extending [方法] to [方向].

% ============================ References ============================
\bibliographystyle{IEEEbib}
\bibliography{refs}   % 需 refs.bib

\end{document}
